\documentclass[11pt]{article}

\usepackage[utf8]{inputenc}
\usepackage[T1]{fontenc}
\usepackage[spanish,es-nodecimaldot]{babel}
\usepackage{amsmath,amssymb,mathtools}
\usepackage{geometry}
\usepackage{booktabs}
\usepackage{array}
\usepackage{enumitem}
\usepackage{hyperref}
\usepackage{xcolor}
\usepackage{siunitx}

\geometry{margin=2.5cm}
\hypersetup{
    colorlinks=true,
    linkcolor=blue!50!black,
    urlcolor=blue!50!black
}

\title{Documento Unificado Depurado\\Plataforma Antivibratoria Activa para Rechazo de Perturbaciones}
\author{Versi\'on principal para estudio, sustentaci\'on e informe}
\date{15 de abril de 2026}

\begin{document}

\maketitle
\tableofcontents
\newpage

\section{Objetivo}

En este documento se deja una sola versi\'on del modelo, con la menor cantidad de elementos innecesarios, pero conservando el rigor mec\'anico y la justificaci\'on f\'isica. La idea es usar un modelo que sea:

\begin{itemize}[leftmargin=1.2cm]
\item defendible con leyes de mec\'anica y circuitos;
\item de tercer orden como pide el proyecto;
\item suficientemente claro para obtener espacio de estados, funci\'on de transferencia y estabilidad;
\item lo bastante limpio para no meter t\'erminos no lineales que no est\'en bien justificados.
\end{itemize}

El problema ya no se plantea como seguimiento de referencia, sino como \textbf{rechazo de perturbaciones}. El objetivo es impedir que las vibraciones de la base se transmitan a la plataforma donde est\'a montado un sensor de ultrasonido usado en pruebas de laboratorio para inspecci\'on de papas.

La decisi\'on metodol\'ogica central es esta:

\begin{quote}
\textbf{el modelo principal se toma lineal y de tercer orden, usando masa, resorte lineal, amortiguador viscoso y un actuador electromec\'anico.}
\end{quote}

Con esto:

\begin{itemize}[leftmargin=1.2cm]
\item no es necesario introducir rigidez c\'ubica;
\item no es obligatorio introducir fricci\'on seca;
\item no hace falta forzar una linealizaci\'on del modelo principal, porque el modelo principal ya es lineal.
\item la estrategia de control se plantea como \textbf{IMC-PID} orientada al rechazo de perturbaciones.
\end{itemize}

\section{Contexto del laboratorio}

En el laboratorio se realizan pruebas sobre papas para detectar plagas usando un sensor de ultrasonido. Ese sensor es sensible al movimiento de la estructura donde se monta. Si la mesa o base vibra, la distancia medida por el sensor cambia aunque el objeto observado no se haya desplazado realmente. Eso introduce error en la lectura.

Por tanto, el problema de control se formula as\'i:

\begin{quote}
\textbf{se desea aislar vibracionalmente el sensor para que las perturbaciones mec\'anicas de la base afecten lo menos posible la medici\'on ultras\'onica.}
\end{quote}

En esta primera aproximaci\'on se asume que:

\begin{itemize}[leftmargin=1.2cm]
\item la perturbaci\'on dominante es traslacional en el eje principal de medici\'on del sensor;
\item la papa observada permanece aproximadamente fija respecto al laboratorio durante la ventana de medida;
\item si se reduce el movimiento de la plataforma del sensor, tambi\'en se reduce el error de medici\'on.
\end{itemize}

\section{Sistema fisico}

Se modela una plataforma antivibratoria activa compuesta por:

\begin{itemize}[leftmargin=1.2cm]
\item una base inferior que vibra por una perturbaci\'on externa;
\item una plataforma superior donde se monta el sensor de ultrasonido;
\item un resorte entre base y plataforma;
\item un amortiguador viscoso entre base y plataforma;
\item un actuador electromec\'anico que aplica fuerza entre base y plataforma.
\end{itemize}

El objetivo del control es reducir el movimiento de la plataforma superior frente al movimiento de la base, para mantener estable la posici\'on del sensor durante la inspecci\'on.

\section{Supuestos minimos del modelo}

Los supuestos usados son pocos y expl\'icitos:

\begin{enumerate}[leftmargin=1.2cm]
\item La plataforma se modela como una masa concentrada $m$.
\item El movimiento se estudia en una sola direcci\'on.
\item El resorte trabaja en el rango lineal de Hooke.
\item El amortiguamiento se aproxima como viscoso lineal.
\item El actuador es electromec\'anico lineal en el rango de operaci\'on.
\item Los par\'ametros son constantes.
\item La variable que m\'as afecta la medici\'on es el desplazamiento del sensor en el eje de medida.
\end{enumerate}

Estos supuestos son los que corresponden al modelo cl\'asico de vibraciones mec\'anicas de un grado de libertad, m\'as una din\'amica el\'ectrica para el actuador.

\section{Variables y glosario basico}

\renewcommand{\arraystretch}{1.2}
\begin{center}
\small
\begin{tabular}{>{\raggedright\arraybackslash}p{2.6cm} >{\raggedright\arraybackslash}p{9.6cm} >{\raggedright\arraybackslash}p{2.0cm}}
\toprule
\textbf{Simbolo} & \textbf{Significado} & \textbf{Unidad} \\
\midrule
$t$ & Tiempo. & s \\
$x(t)$ & Posici\'on absoluta de la plataforma. & m \\
$z(t)$ & Posici\'on absoluta de la base. & m \\
$q(t)$ & Desplazamiento relativo entre plataforma y base: $q=x-z$. & m \\
$\dot q(t)$ & Velocidad relativa. & m/s \\
$\ddot q(t)$ & Aceleraci\'on relativa. & m/s$^2$ \\
$\ddot z(t)$ & Aceleraci\'on de la base. Perturbaci\'on externa principal. & m/s$^2$ \\
$u(t)$ & Voltaje de entrada aplicado al actuador. & V \\
$i(t)$ & Corriente del actuador. & A \\
$F_a(t)$ & Fuerza aplicada por el actuador. & N \\
$m$ & Masa equivalente de la plataforma. & kg \\
$k$ & Rigidez lineal del resorte. & N/m \\
$c$ & Coeficiente de amortiguamiento viscoso. & N\,s/m \\
$L$ & Inductancia del actuador. & H \\
$R$ & Resistencia del actuador. & $\Omega$ \\
$K_f$ & Constante de fuerza del actuador, definida por $F_a=K_f i$. & N/A \\
$K_e$ & Constante de back-EMF, definida por $e_b=K_e\dot q$. & V\,s/m \\
$y(t)$ & Salida regulada del sistema. En este problema conviene tomarla como la posici\'on absoluta de la plataforma, $y=x(t)$. & m \\
$h(t)$ & Distancia medida por el sensor ultras\'onico. Si la papa est\'a quieta respecto al laboratorio, las variaciones de $h(t)$ quedan directamente afectadas por el movimiento de la plataforma. & m \\
$r(t)$ & Referencia de regulaci\'on. En rechazo de perturbaciones se toma t\'ipicamente $r(t)=0$. & m \\
$s$ & Variable del dominio de Laplace. & 1/s \\
\bottomrule
\end{tabular}
\end{center}

\subsection{Notacion}

\begin{itemize}[leftmargin=1.2cm]
\item $\dot r$ indica derivada temporal de una variable $r$.
\item $\ddot r$ indica segunda derivada temporal.
\item Las letras may\'usculas como $Q(s)$ o $U(s)$ representan transformadas de Laplace.
\end{itemize}

\section{Planteamiento como problema de rechazo de perturbaciones}

La referencia natural del problema es que la plataforma permanezca quieta. Por eso se plantea:

\begin{equation}
r(t)=0.
\end{equation}

La salida que se desea regular es la posici\'on absoluta de la plataforma donde va montado el sensor:
\begin{equation}
y(t)=x(t).
\end{equation}

La perturbaci\'on f\'isica es la vibraci\'on de la base, representada por $z(t)$ o por su aceleraci\'on $\ddot z(t)$. Sin embargo, hay que hacer una distinci\'on metodol\'ogica importante:

\begin{quote}
\textbf{la perturbaci\'on no se tratar\'a como una entrada adicional de la planta nominal cuando se obtenga la representaci\'on en espacio de estados y la funci\'on de transferencia principal.}
\end{quote}

Es decir:

\begin{itemize}[leftmargin=1.2cm]
\item la \textbf{entrada} del modelo nominal es solo el voltaje de control $u(t)$;
\item la \textbf{perturbaci\'on} es una se\~nal externa que act\'ua sobre el sistema, pero no forma parte del vector de entrada usado para construir la tripleta $A,B,C$ y la funci\'on de transferencia nominal.
\end{itemize}

Entonces el objetivo del control es:
\begin{equation}
y(t)\approx 0 \quad \text{a pesar de la perturbacion de base}.
\end{equation}

En t\'erminos experimentales, esto significa que la mesa puede vibrar, pero la plataforma del sensor debe permanecer lo m\'as inm\'ovil posible para que la medici\'on ultras\'onica no se degrade.

\subsection{Resumen estructural corregido}

\begin{center}
\small
\begin{tabular}{>{\raggedright\arraybackslash}p{3.2cm} >{\raggedright\arraybackslash}p{2.8cm} >{\raggedright\arraybackslash}p{7.4cm}}
\toprule
\textbf{Categoria} & \textbf{Simbolo} & \textbf{Interpretacion fisica} \\
\midrule
Entrada de control & $u(t)$ & Voltaje aplicado al actuador. Es la unica entrada de la planta nominal para obtener $A$, $B$, $C$ y la funcion de transferencia principal. \\
Perturbacion externa & $z(t)$ o $\ddot z(t)$ & Vibracion de la base. Afecta el sistema fisico, pero no se incorpora como una segunda entrada de la planta nominal al derivar la TF principal. \\
Salida regulada & $y(t)=x(t)$ & Posicion absoluta de la plataforma donde esta montado el sensor ultras\'onico. Es la variable que se quiere mantener cercana a cero. \\
Variable interna de modelado & $q(t)=x(t)-z(t)$ & Desplazamiento relativo entre plataforma y base. Facilita escribir las fuerzas del resorte y del amortiguador. \\
Estado 1 & $x_1=q$ & Deformacion relativa de la suspension. \\
Estado 2 & $x_2=\dot q$ & Velocidad relativa entre plataforma y base. \\
Estado 3 & $x_3=i$ & Corriente del actuador. No es una entrada externa, sino una variable interna determinada por la dinamica electrica. \\
Orden del modelo nominal & $3$ & El sistema es de tercer orden porque requiere tres estados independientes: $q$, $\dot q$ e $i$. \\
\bottomrule
\end{tabular}
\end{center}

La conclusi\'on importante de la tabla es que \textbf{la entrada manipulada es el voltaje $u(t)$, no la corriente}. La corriente $i(t)$ es un estado interno del actuador, y precisamente por eso aparece el tercer orden en la planta nominal.

\section{Coordenada relativa}

Se define:
\begin{equation}
q(t)=x(t)-z(t).
\label{eq:qdef}
\end{equation}

Esta es la coordenada correcta porque:

\begin{itemize}[leftmargin=1.2cm]
\item el resorte depende de la deformaci\'on relativa;
\item el amortiguador depende de la velocidad relativa;
\item la vibraci\'on de base entra de manera natural en la ecuaci\'on.
\end{itemize}

Derivando:
\begin{equation}
\dot q(t)=\dot x(t)-\dot z(t),
\end{equation}
\begin{equation}
\ddot q(t)=\ddot x(t)-\ddot z(t).
\end{equation}

\section{Modelo mecanico}

\subsection{Segunda ley de Newton}

Sobre la plataforma se aplica:
\begin{equation}
\sum F = m\ddot x.
\label{eq:newton}
\end{equation}

\subsection{Fuerza del resorte}

Por ley de Hooke:
\begin{equation}
F_s = -k(x-z) = -kq.
\label{eq:hooke}
\end{equation}

\subsection{Fuerza del amortiguador}

Para el amortiguador viscoso:
\begin{equation}
F_d = -c(\dot x-\dot z) = -c\dot q.
\label{eq:damper}
\end{equation}

\subsection{Fuerza del actuador}

La fuerza activa se denota por $F_a(t)$ y act\'ua entre base y plataforma.

\subsection{Ecuacion mecanica en coordenada absoluta}

Sumando fuerzas:
\begin{equation}
m\ddot x = -k(x-z)-c(\dot x-\dot z)+F_a.
\label{eq:mech_abs}
\end{equation}

\subsection{Paso a coordenada relativa}

Usando $x=q+z$:
\begin{equation}
m(\ddot q+\ddot z) = -kq-c\dot q+F_a.
\end{equation}

Reordenando:
\begin{equation}
m\ddot q + c\dot q + kq = F_a - m\ddot z.
\label{eq:mech_rel}
\end{equation}

Esta ecuaci\'on muestra que la aceleraci\'on de base act\'ua como una excitaci\'on inercial equivalente.

\section{Modelo del actuador}

Aqu\'i es donde aparece el tercer orden del sistema.

\subsection{Relacion fuerza-corriente}

Se asume un actuador electromec\'anico lineal para el cual la fuerza es proporcional a la corriente:
\begin{equation}
F_a = K_f i.
\label{eq:force_current}
\end{equation}

Esta ecuaci\'on significa que el actuador convierte corriente el\'ectrica en fuerza mec\'anica.

\subsection{Ecuacion electrica}

Por ley de voltajes de Kirchhoff:
\begin{equation}
u = L\dot i + Ri + e_b.
\label{eq:kvl}
\end{equation}

El t\'ermino $e_b$ es la fuerza contraelectromotriz generada por el movimiento del actuador. Se modela como:
\begin{equation}
e_b = K_e \dot q.
\label{eq:bemf}
\end{equation}

Entonces la ecuaci\'on el\'ectrica completa queda:
\begin{equation}
L\dot i + Ri + K_e\dot q = u.
\label{eq:elec}
\end{equation}

\subsection{Interpretacion fisica del actuador}

La cadena fisica correcta es:

\begin{equation*}
u(t)\longrightarrow i(t)\longrightarrow F_a(t)\longrightarrow q(t).
\end{equation*}

Es decir:

\begin{itemize}[leftmargin=1.2cm]
\item el controlador entrega un voltaje $u(t)$;
\item ese voltaje produce una corriente $i(t)$;
\item la corriente genera una fuerza $F_a(t)$;
\item la fuerza corrige el movimiento de la plataforma.
\end{itemize}

Por eso \textbf{no} es correcto identificar directamente $u$ con fuerza.

\section{Modelo completo exacto}

Sustituyendo \eqref{eq:force_current} en \eqref{eq:mech_rel}:
\begin{equation}
m\ddot q + c\dot q + kq = K_f i - m\ddot z.
\label{eq:full_mech}
\end{equation}

Junto con \eqref{eq:elec}:
\begin{equation}
L\dot i + Ri + K_e\dot q = u.
\label{eq:full_elec}
\end{equation}

Estas dos ecuaciones forman el modelo principal del proyecto.

\subsection{Por que este modelo ya cumple el orden minimo}

El sistema ya es de tercer orden porque tiene tres estados naturales:

\begin{equation}
x_1=q,\qquad x_2=\dot q,\qquad x_3=i.
\end{equation}

No es necesario meter rigidez c\'ubica para obtener orden tres.

\section{Espacio de estados nominal}

Para obtener la representaci\'on nominal en espacio de estados, la perturbaci\'on de base se anula temporalmente:
\[
z(t)=0,\qquad \dot z(t)=0,\qquad \ddot z(t)=0.
\]

En ese caso, la salida absoluta y la relativa coinciden:
\[
x(t)=q(t).
\]

Definiendo los estados:
\[
x_1=q,\qquad x_2=\dot q,\qquad x_3=i,
\]
las ecuaciones quedan:
\begin{equation}
\dot x_1 = x_2,
\end{equation}
\begin{equation}
\dot x_2 = -\frac{k}{m}x_1 - \frac{c}{m}x_2 + \frac{K_f}{m}x_3,
\end{equation}
\begin{equation}
\dot x_3 = -\frac{K_e}{L}x_2 - \frac{R}{L}x_3 + \frac{1}{L}u.
\end{equation}

Por tanto, la planta nominal queda en la forma:
\begin{equation}
\dot x = Ax + Bu,
\label{eq:ss_nom_1}
\end{equation}
\begin{equation}
y = Cx + Du,
\label{eq:ss_nom_2}
\end{equation}
con
\begin{equation}
A=
\begin{bmatrix}
0 & 1 & 0\\
-\dfrac{k}{m} & -\dfrac{c}{m} & \dfrac{K_f}{m}\\
0 & -\dfrac{K_e}{L} & -\dfrac{R}{L}
\end{bmatrix},
\qquad
B=
\begin{bmatrix}
0\\
0\\
\dfrac{1}{L}
\end{bmatrix},
\label{eq:AB_nom}
\end{equation}
\begin{equation}
C=
\begin{bmatrix}
1 & 0 & 0
\end{bmatrix},
\qquad
D=0.
\label{eq:CD_nom}
\end{equation}

Aqu\'i se est\'a usando como salida nominal:
\[
y=q.
\]

Eso es consistente con la derivaci\'on porque, al anular la perturbaci\'on, se cumple que $q=x$.

\section{Funcion de transferencia nominal}

Para obtener la planta nominal de control se toma perturbaci\'on de base nula:
\[
\ddot z = 0.
\]

Entonces:
\begin{equation}
m\ddot q + c\dot q + kq = K_f i,
\label{eq:tf_mech}
\end{equation}
\begin{equation}
L\dot i + Ri + K_e\dot q = u.
\label{eq:tf_elec}
\end{equation}

Aplicando Laplace con condiciones iniciales nulas:
\[
(ms^2+cs+k)Q(s)=K_f I(s),
\]
\[
(Ls+R)I(s)+K_e sQ(s)=U(s).
\]

Despejando $I(s)$:
\[
I(s)=\frac{ms^2+cs+k}{K_f}Q(s).
\]

Sustituyendo en la ecuaci\'on el\'ectrica:
\[
(Ls+R)\frac{ms^2+cs+k}{K_f}Q(s)+K_e sQ(s)=U(s).
\]

Multiplicando por $K_f$:
\[
\left[(Ls+R)(ms^2+cs+k)+K_fK_e s\right]Q(s)=K_fU(s).
\]

Por tanto:
\begin{equation}
G(s)=\frac{Q(s)}{U(s)}
=
\frac{K_f}{(Ls+R)(ms^2+cs+k)+K_fK_e s}.
\label{eq:G}
\end{equation}

Expandiendo:
\begin{equation}
G(s)=
\frac{K_f}
{Lm\,s^3 + (Lc+Rm)s^2 + (Lk+Rc+K_fK_e)s + Rk}.
\label{eq:Gexpanded}
\end{equation}

Esta funci\'on de transferencia describe el canal nominal desde la entrada de control $u(t)$ hasta la salida $y(t)$ cuando la base no se mueve.

\subsection{Obtencion de la TF desde espacio de estados}

Como la planta nominal se escribi\'o en la forma:
\[
\dot x = Ax+Bu,\qquad y=Cx+Du,
\]
la funci\'on de transferencia se obtiene con la expresi\'on:
\begin{equation}
G(s)=C(sI-A)^{-1}B+D.
\label{eq:ss2tf}
\end{equation}

Sustituyendo las matrices \eqref{eq:AB_nom}--\eqref{eq:CD_nom}, se obtiene exactamente la misma expresi\'on de \eqref{eq:Gexpanded}. Esa es la funci\'on de transferencia principal que se debe usar para el dise\~no del controlador.

\subsection{Donde queda la perturbacion}

La perturbaci\'on no desaparece del problema f\'isico; simplemente \textbf{no se incorpora como segunda entrada en la TF nominal}. Su papel correcto es:

\begin{itemize}[leftmargin=1.2cm]
\item en la derivaci\'on del modelo completo aparece en la ecuaci\'on mec\'anica como $-m\ddot z$;
\item en el an\'alisis nominal para obtener $A$, $B$, $C$ y $G(s)$ se fija temporalmente en cero;
\item en el esquema de control se interpreta como una se\~nal externa que perturba la planta y cuyo efecto debe ser rechazado por el lazo cerrado.
\end{itemize}

Esta separaci\'on es la que hace consistente el enfoque de rechazo de perturbaciones con la representaci\'on cl\'asica:
\[
\dot x=Ax+Bu,\qquad y=Cx+Du.
\]

\section{Estrategia de control IMC-PID}

En este proyecto la planta f\'isica principal sigue siendo de \textbf{tercer orden}. Sin embargo, para obtener un controlador de tipo PID mediante IMC, se usa un \textbf{modelo reducido de sinton\'ia}. Esta es una pr\'actica normal en control: la planta completa se conserva para an\'alisis y validaci\'on, mientras que el ajuste del PID se hace sobre una aproximaci\'on m\'as simple pero f\'isicamente coherente.

\subsection{Por que se reduce la planta para sintonizar el PID}

La planta exacta es:
\begin{equation}
m\ddot q + c\dot q + kq = K_f i - m\ddot z,
\end{equation}
\begin{equation}
L\dot i + Ri + K_e\dot q = u.
\end{equation}

Esta planta tiene tres estados:
\[
q,\qquad \dot q,\qquad i.
\]

Por tanto, la planta \textbf{s\'i es de tercer orden}. Pero un PID no se dise\~na normalmente sobre la inversa exacta de una planta c\'ubica si lo que se quiere es una estructura sencilla y justificable en laboratorio. Por eso se toma una reducci\'on de orden para el dise\~no del controlador.

\subsection{Modelo reducido para sintonia}

Se asume que la din\'amica el\'ectrica del actuador es m\'as r\'apida que la din\'amica mec\'anica dominante, de modo que en la banda principal de rechazo de perturbaciones puede aproximarse:
\begin{equation}
L\dot i \approx 0.
\label{eq:current_fast}
\end{equation}

Entonces, a partir de la ecuaci\'on el\'ectrica:
\begin{equation}
Ri + K_e\dot q \approx u,
\end{equation}
y por tanto:
\begin{equation}
i \approx \frac{u-K_e\dot q}{R}.
\label{eq:i_approx}
\end{equation}

Sustituyendo \eqref{eq:i_approx} en la ecuaci\'on mec\'anica:
\begin{equation}
m\ddot q + c\dot q + kq
=
K_f\left(\frac{u-K_e\dot q}{R}\right)-m\ddot z.
\end{equation}

Reordenando:
\begin{equation}
m\ddot q + \left(c+\frac{K_fK_e}{R}\right)\dot q + kq
=
\frac{K_f}{R}u - m\ddot z.
\label{eq:reduced_mech}
\end{equation}

Definimos:
\begin{equation}
c_e = c+\frac{K_fK_e}{R},
\qquad
K_u = \frac{K_f}{R}.
\label{eq:reduced_params}
\end{equation}

Entonces el modelo reducido queda:
\begin{equation}
m\ddot q + c_e\dot q + kq = K_u u - m\ddot z.
\label{eq:reduced_model}
\end{equation}

Tomando inicialmente $\ddot z=0$ para la sinton\'ia nominal:
\begin{equation}
G_r(s)=\frac{Q(s)}{U(s)}
=
\frac{K_u}{ms^2+c_e s+k}.
\label{eq:Greduced}
\end{equation}

Este es un modelo de \textbf{segundo orden}, adecuado para obtener un PID por IMC.

\subsection{Forma deseada del lazo cerrado}

Para rechazo de perturbaciones se escoge una din\'amica nominal simple:
\begin{equation}
T_d(s)=\frac{1}{\lambda s+1},
\label{eq:Td}
\end{equation}
donde $\lambda>0$ es el par\'ametro de sinton\'ia IMC.

La interpretaci\'on de $\lambda$ es:

\begin{itemize}[leftmargin=1.2cm]
\item $\lambda$ peque\~no: rechazo m\'as r\'apido, pero mayor esfuerzo de control;
\item $\lambda$ grande: respuesta m\'as suave y m\'as robusta, pero m\'as lenta.
\end{itemize}

\subsection{Obtencion del controlador IMC-PID}

Usando la relaci\'on de dise\~no directo equivalente a IMC:
\begin{equation}
C(s)=\frac{T_d(s)}{G_r(s)\left[1-T_d(s)\right]},
\label{eq:ds}
\end{equation}
y sustituyendo \eqref{eq:Greduced} y \eqref{eq:Td}:
\begin{equation}
C(s)
=
\frac{\dfrac{1}{\lambda s+1}}
{\dfrac{K_u}{ms^2+c_e s+k}\left(1-\dfrac{1}{\lambda s+1}\right)}.
\end{equation}

Como
\[
1-\frac{1}{\lambda s+1}=\frac{\lambda s}{\lambda s+1},
\]
se simplifica a:
\begin{equation}
C(s)=\frac{ms^2+c_e s+k}{K_u\lambda s}.
\label{eq:Cpid_raw}
\end{equation}

Separando t\'erminos:
\begin{equation}
C(s)
=
\frac{m}{K_u\lambda}s
+
\frac{c_e}{K_u\lambda}
+
\frac{k}{K_u\lambda}\frac{1}{s}.
\label{eq:Cpid_split}
\end{equation}

Esta es exactamente la estructura de un PID ideal:
\begin{equation}
C_{PID}(s)=K_c\left(1+\frac{1}{\tau_I s}+\tau_D s\right).
\label{eq:pid_ideal}
\end{equation}

Comparando \eqref{eq:Cpid_split} con \eqref{eq:pid_ideal}, se obtienen los par\'ametros:
\begin{equation}
K_c = \frac{c_e}{K_u\lambda},
\label{eq:Kc1}
\end{equation}
\begin{equation}
\tau_I = \frac{c_e}{k},
\label{eq:Ti1}
\end{equation}
\begin{equation}
\tau_D = \frac{m}{c_e}.
\label{eq:Td1}
\end{equation}

Sustituyendo \eqref{eq:reduced_params}:
\begin{equation}
K_c = \frac{Rc+K_fK_e}{K_f\lambda},
\label{eq:Kc_final}
\end{equation}
\begin{equation}
\tau_I = \frac{Rc+K_fK_e}{Rk},
\label{eq:Ti_final}
\end{equation}
\begin{equation}
\tau_D = \frac{mR}{Rc+K_fK_e}.
\label{eq:Td_final}
\end{equation}

\subsection{Forma practica del PID}

En implementaci\'on real no conviene usar la derivada ideal pura, porque amplifica ruido. Por eso se recomienda usar una derivada filtrada:
\begin{equation}
C_{PID,f}(s)=
K_c\left(
1+\frac{1}{\tau_I s}+\frac{\tau_D s}{\left(\dfrac{\tau_D}{N}\right)s+1}
\right),
\label{eq:pid_filtered}
\end{equation}
donde $N$ es un factor de filtrado, normalmente entre 10 y 20.

\subsection{Por que este controlador es adecuado para rechazo de perturbaciones}

El IMC-PID es apropiado para este proyecto por cuatro razones:

\begin{itemize}[leftmargin=1.2cm]
\item la acci\'on proporcional aporta correcci\'on inmediata ante desviaciones;
\item la acci\'on integral ayuda a eliminar error estacionario causado por perturbaciones lentas o sesgos;
\item la acci\'on derivativa mejora la anticipaci\'on frente a vibraciones r\'apidas;
\item el par\'ametro $\lambda$ permite regular el compromiso entre rapidez y robustez.
\end{itemize}

\subsection{Lectura correcta del orden del sistema}

Aqu\'i hay una aclaraci\'on importante:

\begin{itemize}[leftmargin=1.2cm]
\item la \textbf{planta f\'isica nominal} sigue siendo de tercer orden;
\item el \textbf{modelo de sinton\'ia IMC-PID} es de segundo orden por reducci\'on f\'isicamente justificada;
\item el controlador PID se dise\~na con el modelo reducido;
\item el desempe\~no final debe evaluarse sobre la planta completa.
\end{itemize}

Por tanto, no hay contradicci\'on entre decir que el sistema es de tercer orden y usar un IMC-PID. Lo que cambia de orden es el \textbf{modelo de dise\~no del controlador}, no la planta f\'isica original.

\section{Estabilidad}

La ecuaci\'on caracter\'istica es:
\[
Lm\,s^3 + (Lc+Rm)s^2 + (Lk+Rc+K_fK_e)s + Rk = 0.
\]

Definiendo:
\[
a_3=Lm,\qquad
a_2=Lc+Rm,\qquad
a_1=Lk+Rc+K_fK_e,\qquad
a_0=Rk,
\]
el criterio de Routh-Hurwitz para un c\'ubico exige:
\begin{equation}
a_3>0,\quad a_2>0,\quad a_1>0,\quad a_0>0,\quad a_2a_1>a_3a_0.
\label{eq:RH}
\end{equation}

Como los par\'ametros fisicos son positivos, la condici\'on decisiva es:
\begin{equation}
(Lc+Rm)(Lk+Rc+K_fK_e) > LmRk.
\label{eq:RHcond}
\end{equation}

\section{Sobre la linealizacion}

En esta versi\'on del proyecto la linealizaci\'on \textbf{no es necesaria para el modelo principal}, porque:

\begin{itemize}[leftmargin=1.2cm]
\item el resorte ya se tom\'o lineal;
\item el amortiguador ya se tom\'o viscoso lineal;
\item la relaci\'on fuerza-corriente del actuador se tom\'o lineal;
\item la ecuaci\'on el\'ectrica tambi\'en es lineal.
\end{itemize}

Entonces el modelo principal ya naci\'o lineal y se puede pasar directamente a:

\begin{itemize}[leftmargin=1.2cm]
\item espacio de estados;
\item funci\'on de transferencia;
\item an\'alisis de estabilidad;
\item dise\~no del controlador.
\end{itemize}

\subsection{Cuando si aparece una linealizacion}

Solo har\'ia falta linealizar si decides a\~nadir al modelo realista alguno de estos efectos:

\begin{itemize}[leftmargin=1.2cm]
\item saturaci\'on del actuador;
\item topes mec\'anicos;
\item fricci\'on seca;
\item rigidez no lineal.
\end{itemize}

Pero esos efectos son opcionales. No son necesarios para que el proyecto tenga orden tres ni para que el modelo base sea fisicamente correcto.

\section{Que se elimino y por que}

En esta versi\'on se eliminan elementos que no son imprescindibles:

\subsection{Rigidez cubica}

Se elimina porque no pertenece al modelo base cl\'asico de vibraciones. Solo se usa si se quiere representar un resorte no lineal tipo Duffing o un elemento elastico con endurecimiento.

\subsection{Friccion seca regularizada}

Se elimina del modelo principal porque complica el an\'alisis y no es necesaria para obtener un sistema de tercer orden. Si se desea, puede dejarse solo para un modelo realista adicional en simulaci\'on.

\subsection{Actuador simplificado de primer orden}

Tampoco se usa como modelo principal, porque ya disponemos de un modelo del actuador m\'as claro fisicamente:
\[
F_a=K_f i,\qquad
L\dot i+Ri+K_e\dot q=u.
\]

Eso hace m\'as f\'acil justificar de donde sale el tercer estado.

\section{Relacion con Simulink o Simscape}

La secuencia correcta del trabajo es:

\begin{enumerate}[leftmargin=1.2cm]
\item definir el sistema fisico;
\item escribir las ecuaciones del modelo base;
\item obtener espacio de estados y funci\'on de transferencia;
\item analizar estabilidad;
\item dise\~nar el controlador;
\item implementar ese modelo en Simulink;
\item si se quiere una versi\'on m\'as realista, agregar saturaci\'on o topes en una segunda simulaci\'on.
\end{enumerate}

Lo importante es no confundir el papel de Simscape:

\begin{quote}
\textbf{Simscape no reemplaza la derivacion matem\'atica. Primero se deriva el modelo; luego Simscape sirve para implementarlo o contrastarlo.}
\end{quote}

\section{Conclusion}

La versi\'on depurada y recomendable del proyecto es:

\begin{equation}
m\ddot q + c\dot q + kq = K_f i - m\ddot z,
\end{equation}
\begin{equation}
L\dot i + Ri + K_e\dot q = u.
\end{equation}

Este modelo:

\begin{itemize}[leftmargin=1.2cm]
\item es mecanicamente defendible;
\item usa solo elementos est\'andar de libros de vibraciones y din\'amica;
\item es de tercer orden;
\item est\'a formulado como un problema de rechazo de perturbaciones;
\item no necesita rigidez c\'ubica;
\item no necesita fricci\'on seca para justificar el proyecto;
\item permite obtener directamente espacio de estados, funci\'on de transferencia y estabilidad;
\item representa de forma coherente el montaje de laboratorio con sensor ultras\'onico sobre una plataforma sensible a vibraciones;
\item permite dise\~nar una estrategia de control IMC-PID sobre un modelo reducido sin perder la interpretaci\'on f\'isica de la planta completa.
\end{itemize}

Si luego quieres una extensi\'on m\'as realista, se puede agregar en simulaci\'on:

\begin{itemize}[leftmargin=1.2cm]
\item saturaci\'on del actuador;
\item topes mec\'anicos.
\end{itemize}

Pero el n\'ucleo del proyecto ya queda completo, coherente y sin elementos innecesarios.

\end{document}
